\documentclass[12pt, a4paper]{article}

% --- BASIC SETTINGS & LANGUAGE ---
\usepackage[utf8]{inputenc}
\usepackage[T1]{fontenc}
\usepackage[czech]{babel}
\usepackage{multicol}

% --- TABLES ---
\usepackage{booktabs}
\usepackage{tabularx}

% --- IMAGE HANDLING ---
\usepackage{graphicx}

% --- HYPERLINKS ---
\usepackage[hidelinks]{hyperref}

% --- DOCUMENT INFORMATION ---
\title{Dokumentace k projektu z předmětu IFJ \\[0.5em] \large VUT FIT \\[0.4em]
\normalsize Varianta: TRP-izp} 
\author{Autoři a rozdělení bodů:\\[0.5em]
\textbf{Jaroslav Renda (xrendaj00) --- 25\%}\\
Jakub Ženata (xzenatj00) --- 25\%\\
Jakub Poláček (xpolacj00) --- 25\%\\
Petr Mácha (xmachap00) --- 25\%\\[1em]
Repozitář: \url{https://github.com/Jaroslav-Renda/IFJ2026}}
\date{\today}

\begin{document}

% Generates the title page
\maketitle

% Generates the table of contents
\newpage
\tableofcontents 


% --- CHAPTERS ---

\section{Rozdělení práce}

Jaroslav Renda -- Vedoucí týmu, syntaktická analýza, precedenční tabulka, návrh LL1 gramatiky, testování, precedenční zásobník
\\Jakub Ženata -- Lexikální analýza, stavový automat, dokumentace, testování
\\Jakub Poláček -- Sémantická analýza, tabulka symbolů, abstraktní syntaktický strom, testování, zásobník, Makefile
\\ Petr Mácha -- Generátor kódu, chybové výstupy, testování, CI


\section{Spolupráce v týmu} 
Spolupráce v týmu probíhala hladce, na začátku semestru proběhla první fyzická schůze kde jsme si domluvili formální záležitosti a rozdělili práci.
Taky jsme se dohodli na vedoucím projektu a občasných schůzkách. Poté již komunikace probíhala online formou přes Discord.
Pro verzování kódu jsme využili Github. Nový kód byl přidáván do vlastní větve, do hlavní větve se dostal jen zkontrolovaný kód.
Pro formátování kódu jsme využili editorconfig a clang-format, názvy proměnných jsou v camelCasu a v celý kód je v angličtině.

\section{Blokové schéma} 

\section{Diagram automatu}

\section{LL tabulka}

\section{Precedenční tabulka}   

\section{Popis jednotlivých modulů}

\section{Využití AI při práci na projektu}  



\end{document}